\section[title={Größen zur Beschreibung einer
Welle},reference={Größen zur Beschreibung einer Welle}]

\startitemize[packed]
\item
  Zentrale Größen zur Beschreibung einer Welle sind ihre Amplitude
  $\hat{y}$, ihre Schwingungsdauer $T$, ihre Frequenz $f$ und ihre
  Phasen- bzw. Ausbreitungsgeschwindigkeit $c$.
\stopitemize

\startbuffer[fig:wellen_beschriftet]
\startMPpage

u := 1cm;

drawarrow (((-1 * pi) - 1) * u,0)--(((2 * pi) + 1) *u,0);

drawarrow (0,-1.5 * u)--(0,1.5 * u);

draw function (1, "x", "sin(x)", (-1 * pi) - 1, (2 * pi) + 1, .01) scaled u;

label.rt(btex $y$ etex, (0, 1.5 * u));
label.top(btex $x$ etex, (((2 * pi) + 1) * u, 0));

draw (-0.1 * u, 1 * u)--(0.1 * u, 1 * u);
draw (-0.1 * u, -1 * u)--(0.1 * u, -1 * u);

label.rt(btex $1$ etex, (0, 1 * u));
label.rt(btex $-1$ etex, (0, -1 * u));

draw (pi * u,0.1 * u)--(pi * u, -0.1 * u);
draw (2pi * u,0.1 * u)--(2pi * u, -0.1 * u);
draw (-1 * pi * u,0.1 * u)--(-1 * pi * u, -0.1 * u);

label.bot(btex $\pi$ etex, (pi * u, 0));
label.bot(btex $2\pi$ etex, (2 * pi * u, 0));
label.bot(btex $-\pi$ etex, (-1 * pi * u, 0));

% ---------- Groessen ---------- %

% ------------ Wellenlaenge ------------ %

% lambda_left := (0, 0.5 * u);
% lambda_right := (2 * pi * u, 0.5 * u);

draw (0, 0.5 * u)--(2 * pi * u, 0.5 * u) withcolor blue;
draw (0, 0.7 * u)--(0, 0.3 * u) withcolor blue;
draw (2 * pi * u, 0.7 * u)--(2 * pi * u, 0.3 * u) withcolor blue;

label.top(btex $\lambda$ etex, (pi * u, 0.5 * u));

% ------------ Amplitude ------------ %

draw (-0.5 * u, 0)--(-0.5 * u,1 * u) withcolor red;
draw (-0.7 * u, 0)--(-0.3 * u, 0) withcolor red;
draw (-0.7 * u, 1 * u)--(-0.3 * u, 1 * u) withcolor red;

label.lft(btex $\hat{y}$ etex, (-0.5 * u, 0.5 * u));

\stopMPpage

\stopbuffer

\startframedtext[width=\dimexpr\textwidth+\rightmarginwidth+\rightmargindistance\relax,frame=off,offset=none] % This spans textwidth + margin width
\placefigure[here][fig:wellen_beschriftet]{Größen anhand der sinus 'Welle'}{
\typesetbuffer[fig:wellen_beschriftet]}
\stopframedtext


\startframedtext[width=\dimexpr\textwidth+\rightmarginwidth+\rightmargindistance\relax,frame=off,offset=none] % This spans textwidth + margin width

% \startplacetable[location=here,title={Größen zum Beschreiben einer Welle, von Leiphi}]
\startplacetable[title={Größen zum Beschreiben einer Welle, von Leiphi}]
% \setupTABLE[each][align=flushleft]
\setupTABLE[r][each][bottomframe=on]
\setupTABLE[r][first][topframe=on]
\setupTABLE[c][first][leftframe=on]
\setupTABLE[c][last][rightframe=on]

\bTABLE[offset=1mm,frame=off]

\bTR \bTD $T$ \eTD 
	\bTD 
	Schwingungsdauer: Zeit, die jedes einzelne Teilchen der harmonischen Welle für
	eine volle Schwingung benötigt. Wir gehen dabei davon aus, dass alle
	schwingenden Teilchen die gleiche Schwingungsdauer wie das erregende Teilchen
	besitzen.
	\eTD 
\eTR

\bTR \bTD $f$ \eTD 
	\bTD 
	(Erreger-)Frequenz: Anzahl der Schwingungsperioden jedes einzelnen Teilchens
	pro Zeiteinheit. Wir gehen dabei davon aus, dass alle schwingenden Teilchen die
	gleiche Frequenz wie das erregende Teilchen besitzen. Es gilt $f=\frac{1}{T}$.
	\eTD
\eTR
\bTR \bTD $\Gamma$ \eTD 
	\bTD 
	Kreisfrequenz: Überstrichener Winkel (im Bogenmaß) jedes einzelne Teilchens pro
	Zeiteinheit. Wir gehen dabei davon aus,  dass alle schwingenden Teilchen die
	gleiche Kreisfrequenz wie das erregende Teilchen besitzen. Es gilt $\Gamma= 2 \cdot \pi \cdot f=\frac{2\cdot \pi}{T}$
	\eTD 
\eTR

\bTR 
	\bTD [nc=2,rightframe=on]
	Bemerkung: Unter den von uns genannten Annahmen sind Amplitude $\hat{y}$ und
	Schwingungsdauer $T$ (bzw. Erregerfrequenz $f$ und Kreisfrequenz $\Gamma$) einer Welle
	allein durch den Erreger der Welle bestimmt.
	\eTD 
\eTR
\bTR \bTD $c$ \eTD 
	\bTD 
	Phasen- oder Ausbreitungsgeschwindigkeit der Welle: Geschwindigkeit,
	mit der sich die Störung über den Wellenträger ausbreitet. Leicht zu bestimmen
	ist $c$, wenn man einen ausgezeichneten Punkt (z.B.
	den Wellenberg) beobachtet.

	Achtung: Die Phasengeschwindigkeit ist nicht
	mit der Geschwindigkeit der von der Welle erfassten Teilchen zu verwechseln.
	\eTD 
\eTR
\bTR
	\bTD [nc=2,rightframe=on] 
	Bemerkung: Die Phasen- oder Ausbreitungsgeschwindigkeit $c$ einer Welle ist
	allein durch den Wellenträger bestimmt.
	\eTD 
\eTR
\bTR \bTD $\lambda$ \eTD 
	\bTD 
	Wellenlänge: x-Abstand eines Teilchens zum nächsten Teilchen im gleichen
	Schwingungszustand (d.h. die beiden Teilchen müssen die gleiche Auslenkung und
	die gleiche Geschwindigkeit haben.

	Anmerkung: Zu Teilchen mit gleichem
	Schwingungszustand sagt man auch gleichphasig schwingende Teilchen.
	\eTD 
\eTR

\eTABLE
\stopplacetable

\stopframedtext % }}}
